%%%%%%%%%%%%%%%%%%%%%%%%%%%%%%%%%%%%%%%%%%%%%%%%%%%%%%%%%%%%%%%%%%%%%%%%%%%%%%

%\subsection{Proof for near-final configurations}
%\label{sec: polynomial - near final}
%%%%%%%%%%%%%%%%%%%%%%%%%%%%%%%%%%%%%%%%%%%%%%%%%%%%%%%%%%%%%%%%%%%%%%%%%%%%%%%

%\subsection{Proof for near-final configurations}
%\label{sec: polynomial - near final}
%%%%%%%%%%%%%%%%%%%%%%%%%%%%%%%%%%%%%%%%%%%%%%%%%%%%%%%%%%%%%%%%%%%%%%%%%%%%%%%

%\subsection{Proof for near-final configurations}
%\label{sec: polynomial - near final}
%%%%%%%%%%%%%%%%%%%%%%%%%%%%%%%%%%%%%%%%%%%%%%%%%%%%%%%%%%%%%%%%%%%%%%%%%%%%%%%

%\subsection{Proof for near-final configurations}
%\label{sec: polynomial - near final}
%\input{6-2_near-final.tex}

In this sub-section we prove that $E$ can be perfectly mixed with precision $0$
in a polynomial number of mixing operations when $\nofdroplets{E}=n$ is a power of $2$.
This extends easily to near-final configurations.

Assume that $n$ is a power of $2$. To establish our upper bound, we
consider a sequence of mixing operations on $E$, where at each step we mix
farthest-apart same-parity concentrations in $E$. These concentrations
are both in the same set $E_\pi$, but the two droplets obtained from mixing can be
either in $E_\pi$ or in $E_{\bar{\pi}}$. Thus this mixing sequence might change
cardinalities and structure of these sets over time. To obtain our bound, we 
will need to analyze how these sets evolve during the segments of mixing operations
that preserve parity (namely when mixing in each $E_\pi$
produces a concentration in the same $E_\pi$).

More specifically, the proof idea is this: By the choice of the mixing pair,
if the mixed concentrations are far apart (say, they differ at least by $\diameter(E)/3$), 
then the value of $\potential(E)$ will significantly decrease, per Observation~\ref{obs: far-appart mix}(b).
Otherwise, if the mixed concentrations are close to each other, the subsets $E_{even}$ and
$E_{odd}$ must be separated by a gap of at least $\diameter(E)/3$. For such a 
separated configuration, the decreases of $\potential(E)$ may be very small (because
the diameters of $E_{even}$ and $E_{odd}$ may be small compared to $\diameter(E)$).
We show, however, that after at most a polynomial number of steps there will be a
mixing operation in some $E_\pi$ that will produce a concentration with
parity $\barpi$. This guarantees that at the next step the diameter of $E_{\barpi}$ is
large, so in this step $\potential(E)$ will decrease significantly.     

The above intuition is formalized in  the proof of Lemma~\ref{lem: small mixes} below.
We first prove the following auxiliary lemma.  

%%%%%%%%%%%%%%%%%%%%%%%%%%%% 

\begin{lemma}\label{lem: same-parity mixability}     
	Let $\pi\in\braced{even,odd}$. 
	Consider a sequence of farthest-apart mixing operations on $E_\pi\subseteq E$, such
	that each mixing in $E_\pi$ produces a concentration with the same parity $\pi$.
	If this sequence contains $4ns(E)$ mixing operations, then after
	this sequence we have $\nofconcentrations{E_\pi} = 1$.
\end{lemma} 

\begin{proof}
	Recall that $\potential_{E_\pi,0}$ denotes the value of $\potential(E_\pi)$ before any mixing has been performed.
	Let $\nofdroplets{E_\pi} = n_\pi$.
	By Observation~\ref{obs: far-appart constant mix} (with $\gamma =1$), if we
	repeatedly mix $\min(E_\pi)$ with $\max(E_\pi)$, then after at most $2n_\pi$ mixes
	the value of $\potential(E_\pi)$ decreases at least by a factor of $\half$.
	Thus, after at most $2n_\pi\log{\potential_{E_\pi,0}}$ such mixing operations we must have
	$\potential(E_\pi) = 0$; that is $\nofconcentrations{E_\pi} = 1$.
	(More precisely, we achieve $\potential(E_\pi) \le 1$. Since all concentrations in $E_\pi$
	have the same parity, by simple calculation, this implies that in fact $\potential(E_\pi) = 0$.
	It is worth to point out here that this argument is different than our earlier argument for
	$\potential(E)$, since in this case $\average(E_\pi)$ may not be integer.)	
	To complete the proof, note that
	$n_\pi\leq n$ and $\potential_{E_\pi,0}\leq\potential_{E,0}$ 
	and therefore $2n_\pi\log{\potential_{E_\pi,0}} \leq 2n\log{\potential_{E,0}} \leq 4ns(E)$.
\end{proof}

%%%%%%%%%%%%%%%%%%%%%%%%%%%%    

\begin{lemma} \label{lem: small mixes}
	After any sequence of at most $8ns(E)$ farthest-apart same-parity 
	mixing operations, $\potential(E)$ decreases at least by a factor of $1 - 1/18n$.
\end{lemma}   

\begin{proof} 
	The proof follows the idea outlined at the beginning of this section. 
	Let $\delta=\diameter(E)$, be the (initial) diameter of $E$.
	If the farthest-apart same-parity  pair $x,y\in E$ satisfies
	$|x-y|\geq \delta/3$, then by Observation~$\ref{obs: far-appart mix}$,
	mixing $x$ and $y$ decreases $\potential(E)$ at least by a factor of $1-1/18n$, 
	and we are done. 
	
	Otherwise, we have that any two same-parity concentrations differ by at most $\delta/3$.
	Consider $E_{even}$ and $E_{odd}$. By the case assumption, these sets are not empty.
	Assume without loss of generality that $\max(E_{odd}) > \max(E_{even})$.
	We then have that $\min(E_{odd}) > \max(E_{odd}) - \delta/3$,
	$\min(E) = \min(E_{even})$  and $\max(E_{even}) < \min(E_{even}) + \delta/3$.
	Hence, $\min(E_{odd})  \ge \max(E_{even}) + \delta/3$.
	
	Now, from Lemma~$\ref{lem: same-parity mixability}$, 
	we derive that it takes at most $4ns(E)$ farthest-apart mixing operations on $E_{even}$ 
	for either an odd concentration to be produced or for $\nofconcentrations{E_{even}} = 1$ to hold 
	(similarly for $E_{odd}$). 
	Thus, if we repeatedly mix farthest-apart same-parity droplets in $E$,
	eventually, after fewer than $8ns(E)$ such mixing operations,	
	either a mixing in $E_{even}$ will produce an odd concentration or 
	a mixing in $E_{odd}$ will produce an even concentration.
	(This is true because we cannot have both $\nofconcentrations{E_{even}} = 1$ 
	and $\nofconcentrations{E_{odd}} = 1$:
	as mentioned in the proof for Lemma~\ref{lem: mixability for near-final},
	$n$ being a power-of-$2$ guarantees the existence of two distinct concentrations with same-parity.)
	
	So, if an odd $x$ was produced from a mixing in $E_{even}$, then we can mix $x$ with $\max(E_{odd})$.
	Otherwise, an even $y$ was produced from a mixing in $E_{odd}$ and we can mix $y$ with $\min(E_{even})$.
	As both $|x-\max(E_{odd})|$ and $|y-\min(E_{even})|$ are at least $\delta/3$, 
	either mixing decreases $\potential(E)$ at least by a factor of $1-1/18n$
	(see Observation~$\ref{obs: far-appart mix}$), and thus the lemma holds.
\end{proof}         

%%%%%%%%%%%%%%%%%%%%%%%%%%%%

Using Lemma~$\ref{lem: small mixes}$ we can now establish a bound on the length of mixing sequences
when $n$ is a power of $2$, and, more generally, when $E$ is near-final.

\begin{theorem} \label{them: n power of two}
	If $\nofdroplets{E}=n$ is a power of $2$, then
	$E$ can be perfectly mixed with precision $0$ 
	by a mixing sequence of length at most $288n^2s^2(E)$.
\end{theorem}

\begin{proof}
	By Lemma~\ref{lem: small mixes}, after a mixing sequence of at most $8ns(E)$ 
	same-parity farthest-apart mixing operations,
	$\potential(E)$ decreases at least by a factor of $ (1-1/18n)$.
	It follows that  after at most $18n$ such mixing sequences, $\potential(E)$ decreases
	at least by a factor of $\half$ (see the proof of Observation~\ref{obs: far-appart constant mix}).   
	Consequently, after at most $18n\log{\potential_{E,0}}$ such mixing sequences,
	$E$ becomes perfectly mixed.	
	Finally, as $\log {\potential_{E,0}} \le 2 s(E)$ and 
	each mixing sequence contains at most $8ns(E)$ mixing operations,
	we obtain that the total number of mixing operations is at most 
	$(18n\log{\potential_{E,0}})\cdot (8ns(E)) \le (36ns(E)) \cdot (8ns(E)) = 288n^2s^2(E)$.
\end{proof}

%%%%%%%%%%%%%%%%%%%%%%%%%%%%   

\begin{theorem}\label{thm: E near-final}
	If $E$ is near-final, then $E$ can be perfectly mixed, with precision $0$,
	by a mixing sequence of length at most $144n^3s^2(E)$.
\end{theorem}

\begin{proof}
	Assume that $E$ is near-final. Recall that $\average(E) = \hatmu$.
	By definition, $E$ can be partitioned into singletons $\hatmu$ and
	at most $n/2$ disjoint multisets $A\subseteq E$, each satisfying $\average(A) = \hatmu$
	and having cardinality that is a (non-zero) power of $2$.
	By Theorem~$\ref{them: n power of two}$ above,
	each multiset $A$ in this partition can be perfectly mixed with at most $288n^2s^2(E)$ mixing operations,
	and the theorem holds.
\end{proof}

In this sub-section we prove that $E$ can be perfectly mixed with precision $0$
in a polynomial number of mixing operations when $\nofdroplets{E}=n$ is a power of $2$.
This extends easily to near-final configurations.

Assume that $n$ is a power of $2$. To establish our upper bound, we
consider a sequence of mixing operations on $E$, where at each step we mix
farthest-apart same-parity concentrations in $E$. These concentrations
are both in the same set $E_\pi$, but the two droplets obtained from mixing can be
either in $E_\pi$ or in $E_{\bar{\pi}}$. Thus this mixing sequence might change
cardinalities and structure of these sets over time. To obtain our bound, we 
will need to analyze how these sets evolve during the segments of mixing operations
that preserve parity (namely when mixing in each $E_\pi$
produces a concentration in the same $E_\pi$).

More specifically, the proof idea is this: By the choice of the mixing pair,
if the mixed concentrations are far apart (say, they differ at least by $\diameter(E)/3$), 
then the value of $\potential(E)$ will significantly decrease, per Observation~\ref{obs: far-appart mix}(b).
Otherwise, if the mixed concentrations are close to each other, the subsets $E_{even}$ and
$E_{odd}$ must be separated by a gap of at least $\diameter(E)/3$. For such a 
separated configuration, the decreases of $\potential(E)$ may be very small (because
the diameters of $E_{even}$ and $E_{odd}$ may be small compared to $\diameter(E)$).
We show, however, that after at most a polynomial number of steps there will be a
mixing operation in some $E_\pi$ that will produce a concentration with
parity $\barpi$. This guarantees that at the next step the diameter of $E_{\barpi}$ is
large, so in this step $\potential(E)$ will decrease significantly.     

The above intuition is formalized in  the proof of Lemma~\ref{lem: small mixes} below.
We first prove the following auxiliary lemma.  

%%%%%%%%%%%%%%%%%%%%%%%%%%%% 

\begin{lemma}\label{lem: same-parity mixability}     
	Let $\pi\in\braced{even,odd}$. 
	Consider a sequence of farthest-apart mixing operations on $E_\pi\subseteq E$, such
	that each mixing in $E_\pi$ produces a concentration with the same parity $\pi$.
	If this sequence contains $4ns(E)$ mixing operations, then after
	this sequence we have $\nofconcentrations{E_\pi} = 1$.
\end{lemma} 

\begin{proof}
	Recall that $\potential_{E_\pi,0}$ denotes the value of $\potential(E_\pi)$ before any mixing has been performed.
	Let $\nofdroplets{E_\pi} = n_\pi$.
	By Observation~\ref{obs: far-appart constant mix} (with $\gamma =1$), if we
	repeatedly mix $\min(E_\pi)$ with $\max(E_\pi)$, then after at most $2n_\pi$ mixes
	the value of $\potential(E_\pi)$ decreases at least by a factor of $\half$.
	Thus, after at most $2n_\pi\log{\potential_{E_\pi,0}}$ such mixing operations we must have
	$\potential(E_\pi) = 0$; that is $\nofconcentrations{E_\pi} = 1$.
	(More precisely, we achieve $\potential(E_\pi) \le 1$. Since all concentrations in $E_\pi$
	have the same parity, by simple calculation, this implies that in fact $\potential(E_\pi) = 0$.
	It is worth to point out here that this argument is different than our earlier argument for
	$\potential(E)$, since in this case $\average(E_\pi)$ may not be integer.)	
	To complete the proof, note that
	$n_\pi\leq n$ and $\potential_{E_\pi,0}\leq\potential_{E,0}$ 
	and therefore $2n_\pi\log{\potential_{E_\pi,0}} \leq 2n\log{\potential_{E,0}} \leq 4ns(E)$.
\end{proof}

%%%%%%%%%%%%%%%%%%%%%%%%%%%%    

\begin{lemma} \label{lem: small mixes}
	After any sequence of at most $8ns(E)$ farthest-apart same-parity 
	mixing operations, $\potential(E)$ decreases at least by a factor of $1 - 1/18n$.
\end{lemma}   

\begin{proof} 
	The proof follows the idea outlined at the beginning of this section. 
	Let $\delta=\diameter(E)$, be the (initial) diameter of $E$.
	If the farthest-apart same-parity  pair $x,y\in E$ satisfies
	$|x-y|\geq \delta/3$, then by Observation~$\ref{obs: far-appart mix}$,
	mixing $x$ and $y$ decreases $\potential(E)$ at least by a factor of $1-1/18n$, 
	and we are done. 
	
	Otherwise, we have that any two same-parity concentrations differ by at most $\delta/3$.
	Consider $E_{even}$ and $E_{odd}$. By the case assumption, these sets are not empty.
	Assume without loss of generality that $\max(E_{odd}) > \max(E_{even})$.
	We then have that $\min(E_{odd}) > \max(E_{odd}) - \delta/3$,
	$\min(E) = \min(E_{even})$  and $\max(E_{even}) < \min(E_{even}) + \delta/3$.
	Hence, $\min(E_{odd})  \ge \max(E_{even}) + \delta/3$.
	
	Now, from Lemma~$\ref{lem: same-parity mixability}$, 
	we derive that it takes at most $4ns(E)$ farthest-apart mixing operations on $E_{even}$ 
	for either an odd concentration to be produced or for $\nofconcentrations{E_{even}} = 1$ to hold 
	(similarly for $E_{odd}$). 
	Thus, if we repeatedly mix farthest-apart same-parity droplets in $E$,
	eventually, after fewer than $8ns(E)$ such mixing operations,	
	either a mixing in $E_{even}$ will produce an odd concentration or 
	a mixing in $E_{odd}$ will produce an even concentration.
	(This is true because we cannot have both $\nofconcentrations{E_{even}} = 1$ 
	and $\nofconcentrations{E_{odd}} = 1$:
	as mentioned in the proof for Lemma~\ref{lem: mixability for near-final},
	$n$ being a power-of-$2$ guarantees the existence of two distinct concentrations with same-parity.)
	
	So, if an odd $x$ was produced from a mixing in $E_{even}$, then we can mix $x$ with $\max(E_{odd})$.
	Otherwise, an even $y$ was produced from a mixing in $E_{odd}$ and we can mix $y$ with $\min(E_{even})$.
	As both $|x-\max(E_{odd})|$ and $|y-\min(E_{even})|$ are at least $\delta/3$, 
	either mixing decreases $\potential(E)$ at least by a factor of $1-1/18n$
	(see Observation~$\ref{obs: far-appart mix}$), and thus the lemma holds.
\end{proof}         

%%%%%%%%%%%%%%%%%%%%%%%%%%%%

Using Lemma~$\ref{lem: small mixes}$ we can now establish a bound on the length of mixing sequences
when $n$ is a power of $2$, and, more generally, when $E$ is near-final.

\begin{theorem} \label{them: n power of two}
	If $\nofdroplets{E}=n$ is a power of $2$, then
	$E$ can be perfectly mixed with precision $0$ 
	by a mixing sequence of length at most $288n^2s^2(E)$.
\end{theorem}

\begin{proof}
	By Lemma~\ref{lem: small mixes}, after a mixing sequence of at most $8ns(E)$ 
	same-parity farthest-apart mixing operations,
	$\potential(E)$ decreases at least by a factor of $ (1-1/18n)$.
	It follows that  after at most $18n$ such mixing sequences, $\potential(E)$ decreases
	at least by a factor of $\half$ (see the proof of Observation~\ref{obs: far-appart constant mix}).   
	Consequently, after at most $18n\log{\potential_{E,0}}$ such mixing sequences,
	$E$ becomes perfectly mixed.	
	Finally, as $\log {\potential_{E,0}} \le 2 s(E)$ and 
	each mixing sequence contains at most $8ns(E)$ mixing operations,
	we obtain that the total number of mixing operations is at most 
	$(18n\log{\potential_{E,0}})\cdot (8ns(E)) \le (36ns(E)) \cdot (8ns(E)) = 288n^2s^2(E)$.
\end{proof}

%%%%%%%%%%%%%%%%%%%%%%%%%%%%   

\begin{theorem}\label{thm: E near-final}
	If $E$ is near-final, then $E$ can be perfectly mixed, with precision $0$,
	by a mixing sequence of length at most $144n^3s^2(E)$.
\end{theorem}

\begin{proof}
	Assume that $E$ is near-final. Recall that $\average(E) = \hatmu$.
	By definition, $E$ can be partitioned into singletons $\hatmu$ and
	at most $n/2$ disjoint multisets $A\subseteq E$, each satisfying $\average(A) = \hatmu$
	and having cardinality that is a (non-zero) power of $2$.
	By Theorem~$\ref{them: n power of two}$ above,
	each multiset $A$ in this partition can be perfectly mixed with at most $288n^2s^2(E)$ mixing operations,
	and the theorem holds.
\end{proof}

In this sub-section we prove that $E$ can be perfectly mixed with precision $0$
in a polynomial number of mixing operations when $\nofdroplets{E}=n$ is a power of $2$.
This extends easily to near-final configurations.

Assume that $n$ is a power of $2$. To establish our upper bound, we
consider a sequence of mixing operations on $E$, where at each step we mix
farthest-apart same-parity concentrations in $E$. These concentrations
are both in the same set $E_\pi$, but the two droplets obtained from mixing can be
either in $E_\pi$ or in $E_{\bar{\pi}}$. Thus this mixing sequence might change
cardinalities and structure of these sets over time. To obtain our bound, we 
will need to analyze how these sets evolve during the segments of mixing operations
that preserve parity (namely when mixing in each $E_\pi$
produces a concentration in the same $E_\pi$).

More specifically, the proof idea is this: By the choice of the mixing pair,
if the mixed concentrations are far apart (say, they differ at least by $\diameter(E)/3$), 
then the value of $\potential(E)$ will significantly decrease, per Observation~\ref{obs: far-appart mix}(b).
Otherwise, if the mixed concentrations are close to each other, the subsets $E_{even}$ and
$E_{odd}$ must be separated by a gap of at least $\diameter(E)/3$. For such a 
separated configuration, the decreases of $\potential(E)$ may be very small (because
the diameters of $E_{even}$ and $E_{odd}$ may be small compared to $\diameter(E)$).
We show, however, that after at most a polynomial number of steps there will be a
mixing operation in some $E_\pi$ that will produce a concentration with
parity $\barpi$. This guarantees that at the next step the diameter of $E_{\barpi}$ is
large, so in this step $\potential(E)$ will decrease significantly.     

The above intuition is formalized in  the proof of Lemma~\ref{lem: small mixes} below.
We first prove the following auxiliary lemma.  

%%%%%%%%%%%%%%%%%%%%%%%%%%%% 

\begin{lemma}\label{lem: same-parity mixability}     
	Let $\pi\in\braced{even,odd}$. 
	Consider a sequence of farthest-apart mixing operations on $E_\pi\subseteq E$, such
	that each mixing in $E_\pi$ produces a concentration with the same parity $\pi$.
	If this sequence contains $4ns(E)$ mixing operations, then after
	this sequence we have $\nofconcentrations{E_\pi} = 1$.
\end{lemma} 

\begin{proof}
	Recall that $\potential_{E_\pi,0}$ denotes the value of $\potential(E_\pi)$ before any mixing has been performed.
	Let $\nofdroplets{E_\pi} = n_\pi$.
	By Observation~\ref{obs: far-appart constant mix} (with $\gamma =1$), if we
	repeatedly mix $\min(E_\pi)$ with $\max(E_\pi)$, then after at most $2n_\pi$ mixes
	the value of $\potential(E_\pi)$ decreases at least by a factor of $\half$.
	Thus, after at most $2n_\pi\log{\potential_{E_\pi,0}}$ such mixing operations we must have
	$\potential(E_\pi) = 0$; that is $\nofconcentrations{E_\pi} = 1$.
	(More precisely, we achieve $\potential(E_\pi) \le 1$. Since all concentrations in $E_\pi$
	have the same parity, by simple calculation, this implies that in fact $\potential(E_\pi) = 0$.
	It is worth to point out here that this argument is different than our earlier argument for
	$\potential(E)$, since in this case $\average(E_\pi)$ may not be integer.)	
	To complete the proof, note that
	$n_\pi\leq n$ and $\potential_{E_\pi,0}\leq\potential_{E,0}$ 
	and therefore $2n_\pi\log{\potential_{E_\pi,0}} \leq 2n\log{\potential_{E,0}} \leq 4ns(E)$.
\end{proof}

%%%%%%%%%%%%%%%%%%%%%%%%%%%%    

\begin{lemma} \label{lem: small mixes}
	After any sequence of at most $8ns(E)$ farthest-apart same-parity 
	mixing operations, $\potential(E)$ decreases at least by a factor of $1 - 1/18n$.
\end{lemma}   

\begin{proof} 
	The proof follows the idea outlined at the beginning of this section. 
	Let $\delta=\diameter(E)$, be the (initial) diameter of $E$.
	If the farthest-apart same-parity  pair $x,y\in E$ satisfies
	$|x-y|\geq \delta/3$, then by Observation~$\ref{obs: far-appart mix}$,
	mixing $x$ and $y$ decreases $\potential(E)$ at least by a factor of $1-1/18n$, 
	and we are done. 
	
	Otherwise, we have that any two same-parity concentrations differ by at most $\delta/3$.
	Consider $E_{even}$ and $E_{odd}$. By the case assumption, these sets are not empty.
	Assume without loss of generality that $\max(E_{odd}) > \max(E_{even})$.
	We then have that $\min(E_{odd}) > \max(E_{odd}) - \delta/3$,
	$\min(E) = \min(E_{even})$  and $\max(E_{even}) < \min(E_{even}) + \delta/3$.
	Hence, $\min(E_{odd})  \ge \max(E_{even}) + \delta/3$.
	
	Now, from Lemma~$\ref{lem: same-parity mixability}$, 
	we derive that it takes at most $4ns(E)$ farthest-apart mixing operations on $E_{even}$ 
	for either an odd concentration to be produced or for $\nofconcentrations{E_{even}} = 1$ to hold 
	(similarly for $E_{odd}$). 
	Thus, if we repeatedly mix farthest-apart same-parity droplets in $E$,
	eventually, after fewer than $8ns(E)$ such mixing operations,	
	either a mixing in $E_{even}$ will produce an odd concentration or 
	a mixing in $E_{odd}$ will produce an even concentration.
	(This is true because we cannot have both $\nofconcentrations{E_{even}} = 1$ 
	and $\nofconcentrations{E_{odd}} = 1$:
	as mentioned in the proof for Lemma~\ref{lem: mixability for near-final},
	$n$ being a power-of-$2$ guarantees the existence of two distinct concentrations with same-parity.)
	
	So, if an odd $x$ was produced from a mixing in $E_{even}$, then we can mix $x$ with $\max(E_{odd})$.
	Otherwise, an even $y$ was produced from a mixing in $E_{odd}$ and we can mix $y$ with $\min(E_{even})$.
	As both $|x-\max(E_{odd})|$ and $|y-\min(E_{even})|$ are at least $\delta/3$, 
	either mixing decreases $\potential(E)$ at least by a factor of $1-1/18n$
	(see Observation~$\ref{obs: far-appart mix}$), and thus the lemma holds.
\end{proof}         

%%%%%%%%%%%%%%%%%%%%%%%%%%%%

Using Lemma~$\ref{lem: small mixes}$ we can now establish a bound on the length of mixing sequences
when $n$ is a power of $2$, and, more generally, when $E$ is near-final.

\begin{theorem} \label{them: n power of two}
	If $\nofdroplets{E}=n$ is a power of $2$, then
	$E$ can be perfectly mixed with precision $0$ 
	by a mixing sequence of length at most $288n^2s^2(E)$.
\end{theorem}

\begin{proof}
	By Lemma~\ref{lem: small mixes}, after a mixing sequence of at most $8ns(E)$ 
	same-parity farthest-apart mixing operations,
	$\potential(E)$ decreases at least by a factor of $ (1-1/18n)$.
	It follows that  after at most $18n$ such mixing sequences, $\potential(E)$ decreases
	at least by a factor of $\half$ (see the proof of Observation~\ref{obs: far-appart constant mix}).   
	Consequently, after at most $18n\log{\potential_{E,0}}$ such mixing sequences,
	$E$ becomes perfectly mixed.	
	Finally, as $\log {\potential_{E,0}} \le 2 s(E)$ and 
	each mixing sequence contains at most $8ns(E)$ mixing operations,
	we obtain that the total number of mixing operations is at most 
	$(18n\log{\potential_{E,0}})\cdot (8ns(E)) \le (36ns(E)) \cdot (8ns(E)) = 288n^2s^2(E)$.
\end{proof}

%%%%%%%%%%%%%%%%%%%%%%%%%%%%   

\begin{theorem}\label{thm: E near-final}
	If $E$ is near-final, then $E$ can be perfectly mixed, with precision $0$,
	by a mixing sequence of length at most $144n^3s^2(E)$.
\end{theorem}

\begin{proof}
	Assume that $E$ is near-final. Recall that $\average(E) = \hatmu$.
	By definition, $E$ can be partitioned into singletons $\hatmu$ and
	at most $n/2$ disjoint multisets $A\subseteq E$, each satisfying $\average(A) = \hatmu$
	and having cardinality that is a (non-zero) power of $2$.
	By Theorem~$\ref{them: n power of two}$ above,
	each multiset $A$ in this partition can be perfectly mixed with at most $288n^2s^2(E)$ mixing operations,
	and the theorem holds.
\end{proof}

In this sub-section we prove that $E$ can be perfectly mixed with precision $0$
in a polynomial number of mixing operations when $\nofdroplets{E}=n$ is a power of $2$.
This extends easily to near-final configurations.

Assume that $n$ is a power of $2$. To establish our upper bound, we
consider a sequence of mixing operations on $E$, where at each step we mix
farthest-apart same-parity concentrations in $E$. These concentrations
are both in the same set $E_\pi$, but the two droplets obtained from mixing can be
either in $E_\pi$ or in $E_{\bar{\pi}}$. Thus this mixing sequence might change
cardinalities and structure of these sets over time. To obtain our bound, we 
will need to analyze how these sets evolve during the segments of mixing operations
that preserve parity (namely when mixing in each $E_\pi$
produces a concentration in the same $E_\pi$).

More specifically, the proof idea is this: By the choice of the mixing pair,
if the mixed concentrations are far apart (say, they differ at least by $\diameter(E)/3$), 
then the value of $\potential(E)$ will significantly decrease, per Observation~\ref{obs: far-appart mix}(b).
Otherwise, if the mixed concentrations are close to each other, the subsets $E_{even}$ and
$E_{odd}$ must be separated by a gap of at least $\diameter(E)/3$. For such a 
separated configuration, the decreases of $\potential(E)$ may be very small (because
the diameters of $E_{even}$ and $E_{odd}$ may be small compared to $\diameter(E)$).
We show, however, that after at most a polynomial number of steps there will be a
mixing operation in some $E_\pi$ that will produce a concentration with
parity $\barpi$. This guarantees that at the next step the diameter of $E_{\barpi}$ is
large, so in this step $\potential(E)$ will decrease significantly.     

The above intuition is formalized in  the proof of Lemma~\ref{lem: small mixes} below.
We first prove the following auxiliary lemma.  

%%%%%%%%%%%%%%%%%%%%%%%%%%%% 

\begin{lemma}\label{lem: same-parity mixability}     
	Let $\pi\in\braced{even,odd}$. 
	Consider a sequence of farthest-apart mixing operations on $E_\pi\subseteq E$, such
	that each mixing in $E_\pi$ produces a concentration with the same parity $\pi$.
	If this sequence contains $4ns(E)$ mixing operations, then after
	this sequence we have $\nofconcentrations{E_\pi} = 1$.
\end{lemma} 

\begin{proof}
	Recall that $\potential_{E_\pi,0}$ denotes the value of $\potential(E_\pi)$ before any mixing has been performed.
	Let $\nofdroplets{E_\pi} = n_\pi$.
	By Observation~\ref{obs: far-appart constant mix} (with $\gamma =1$), if we
	repeatedly mix $\min(E_\pi)$ with $\max(E_\pi)$, then after at most $2n_\pi$ mixes
	the value of $\potential(E_\pi)$ decreases at least by a factor of $\half$.
	Thus, after at most $2n_\pi\log{\potential_{E_\pi,0}}$ such mixing operations we must have
	$\potential(E_\pi) = 0$; that is $\nofconcentrations{E_\pi} = 1$.
	(More precisely, we achieve $\potential(E_\pi) \le 1$. Since all concentrations in $E_\pi$
	have the same parity, by simple calculation, this implies that in fact $\potential(E_\pi) = 0$.
	It is worth to point out here that this argument is different than our earlier argument for
	$\potential(E)$, since in this case $\average(E_\pi)$ may not be integer.)	
	To complete the proof, note that
	$n_\pi\leq n$ and $\potential_{E_\pi,0}\leq\potential_{E,0}$ 
	and therefore $2n_\pi\log{\potential_{E_\pi,0}} \leq 2n\log{\potential_{E,0}} \leq 4ns(E)$.
\end{proof}

%%%%%%%%%%%%%%%%%%%%%%%%%%%%    

\begin{lemma} \label{lem: small mixes}
	After any sequence of at most $8ns(E)$ farthest-apart same-parity 
	mixing operations, $\potential(E)$ decreases at least by a factor of $1 - 1/18n$.
\end{lemma}   

\begin{proof} 
	The proof follows the idea outlined at the beginning of this section. 
	Let $\delta=\diameter(E)$, be the (initial) diameter of $E$.
	If the farthest-apart same-parity  pair $x,y\in E$ satisfies
	$|x-y|\geq \delta/3$, then by Observation~$\ref{obs: far-appart mix}$,
	mixing $x$ and $y$ decreases $\potential(E)$ at least by a factor of $1-1/18n$, 
	and we are done. 
	
	Otherwise, we have that any two same-parity concentrations differ by at most $\delta/3$.
	Consider $E_{even}$ and $E_{odd}$. By the case assumption, these sets are not empty.
	Assume without loss of generality that $\max(E_{odd}) > \max(E_{even})$.
	We then have that $\min(E_{odd}) > \max(E_{odd}) - \delta/3$,
	$\min(E) = \min(E_{even})$  and $\max(E_{even}) < \min(E_{even}) + \delta/3$.
	Hence, $\min(E_{odd})  \ge \max(E_{even}) + \delta/3$.
	
	Now, from Lemma~$\ref{lem: same-parity mixability}$, 
	we derive that it takes at most $4ns(E)$ farthest-apart mixing operations on $E_{even}$ 
	for either an odd concentration to be produced or for $\nofconcentrations{E_{even}} = 1$ to hold 
	(similarly for $E_{odd}$). 
	Thus, if we repeatedly mix farthest-apart same-parity droplets in $E$,
	eventually, after fewer than $8ns(E)$ such mixing operations,	
	either a mixing in $E_{even}$ will produce an odd concentration or 
	a mixing in $E_{odd}$ will produce an even concentration.
	(This is true because we cannot have both $\nofconcentrations{E_{even}} = 1$ 
	and $\nofconcentrations{E_{odd}} = 1$:
	as mentioned in the proof for Lemma~\ref{lem: mixability for near-final},
	$n$ being a power-of-$2$ guarantees the existence of two distinct concentrations with same-parity.)
	
	So, if an odd $x$ was produced from a mixing in $E_{even}$, then we can mix $x$ with $\max(E_{odd})$.
	Otherwise, an even $y$ was produced from a mixing in $E_{odd}$ and we can mix $y$ with $\min(E_{even})$.
	As both $|x-\max(E_{odd})|$ and $|y-\min(E_{even})|$ are at least $\delta/3$, 
	either mixing decreases $\potential(E)$ at least by a factor of $1-1/18n$
	(see Observation~$\ref{obs: far-appart mix}$), and thus the lemma holds.
\end{proof}         

%%%%%%%%%%%%%%%%%%%%%%%%%%%%

Using Lemma~$\ref{lem: small mixes}$ we can now establish a bound on the length of mixing sequences
when $n$ is a power of $2$, and, more generally, when $E$ is near-final.

\begin{theorem} \label{them: n power of two}
	If $\nofdroplets{E}=n$ is a power of $2$, then
	$E$ can be perfectly mixed with precision $0$ 
	by a mixing sequence of length at most $288n^2s^2(E)$.
\end{theorem}

\begin{proof}
	By Lemma~\ref{lem: small mixes}, after a mixing sequence of at most $8ns(E)$ 
	same-parity farthest-apart mixing operations,
	$\potential(E)$ decreases at least by a factor of $ (1-1/18n)$.
	It follows that  after at most $18n$ such mixing sequences, $\potential(E)$ decreases
	at least by a factor of $\half$ (see the proof of Observation~\ref{obs: far-appart constant mix}).   
	Consequently, after at most $18n\log{\potential_{E,0}}$ such mixing sequences,
	$E$ becomes perfectly mixed.	
	Finally, as $\log {\potential_{E,0}} \le 2 s(E)$ and 
	each mixing sequence contains at most $8ns(E)$ mixing operations,
	we obtain that the total number of mixing operations is at most 
	$(18n\log{\potential_{E,0}})\cdot (8ns(E)) \le (36ns(E)) \cdot (8ns(E)) = 288n^2s^2(E)$.
\end{proof}

%%%%%%%%%%%%%%%%%%%%%%%%%%%%   

\begin{theorem}\label{thm: E near-final}
	If $E$ is near-final, then $E$ can be perfectly mixed, with precision $0$,
	by a mixing sequence of length at most $144n^3s^2(E)$.
\end{theorem}

\begin{proof}
	Assume that $E$ is near-final. Recall that $\average(E) = \hatmu$.
	By definition, $E$ can be partitioned into singletons $\hatmu$ and
	at most $n/2$ disjoint multisets $A\subseteq E$, each satisfying $\average(A) = \hatmu$
	and having cardinality that is a (non-zero) power of $2$.
	By Theorem~$\ref{them: n power of two}$ above,
	each multiset $A$ in this partition can be perfectly mixed with at most $288n^2s^2(E)$ mixing operations,
	and the theorem holds.
\end{proof}